\paragraph{Comparison with compact constructions.}\label{sec:benchcomparison}
\Cref{tab:eval-aux} gives separate observations of upstream SHRINCS C B32~\cite{shrincsc} and two compact implementations motivated by Kudinov--Nick~\cite{kn}. The latter are our C reconstruction and third-party Go code~\cite{hashsiggo}, rather than the authors' optimized implementation; we have not established a security reduction for their exact encodings. The different workload, state and preprocessing assumptions preclude pooling these timings with \Cref{tab:eval-warm}.

\begin{table}[t]
\caption{Auxiliary observations on 32-byte messages under background load. Latencies are milliseconds; $N$ counts calls per operation. The security and state assumptions of these implementations are not normalized.}\label{tab:eval-aux}
\centering\small\setlength{\tabcolsep}{3pt}
\begin{tabular}{lrrrr}\toprule
Implementation & Signature bytes & $N$ & Sign median & Verify median\\\midrule
SHRINCS C B32, stateful &324--3,508&200&48.520&0.120\\
SHRINCS C B32, forced fallback &3,680&20&471.961&0.448\\
Compact C reconstruction &3,952&20&3613.611&0.542\\
Compact third-party Go &4,036&12&1897.784&0.905\\\bottomrule
\end{tabular}
\smallskip\par\raggedright\footnotesize The stateful size range covers the 200 measured indices, not a certified signing capacity. The compact configurations refer to published $2^{40}$-query examples; that reference does not certify these implementations. Forced fallback is neither measured reliable loss detection nor a naturally exhausted stateful history.
\end{table}

Silent reuse of a consumption index and honest restoration of LULA I's cache have different security meanings. The target-sum reuse argument is a model analysis; no forgery was run against a real SHRINCS implementation. Neither table establishes a security-normalized ranking.
